\scanpage{10}\mspage{\handwritten{17}}
$H^{3}(Y,\underline G_{\mathrm m})\handwritten{\simeq}\underline{\mathrm Q}/\underline{\mathrm Z}$, hence the corollary.

Let \deleted{\unsure{$i$:}}
\[
  i:\handwritten{\eta\longrightarrow}Y
\]
be the inclusion of the generic point, and put
\[
  \underline B=i_{\handwritten{\ast}}(i^{\handwritten{\ast}}(P)) \quad , \tag*{(4.8)}
\]
so that we have a canonical homomorphism
\[
  P\handwritten{\longrightarrow}\underline B \quad , \tag*{(4.9)}
\]
and we shall put
\[
  \underline E=\operatorname{Ker}(P\handwritten{\to}\underline B)
  \quad ,\quad \underline F=\operatorname{Coker}(P\handwritten{\to}\underline B)\quad . \tag*{(4.10)}
\]

The sheaves $\underline E$ and $\underline F$ are ``skyscraper sheaves'' , which are therefore known when one knows their restrictions to the closed points of $Y$. The sheave\handwritten{s}' geometric fiber\handwritten{s} over such a point $y$ are obtained as the kernels and cokernel of the homomorphism
\[
  \deleted{\mathrm{Pic}(}\tikz[baseline=0pt,x=1em,y=1ex,line cap=round]{
  \path[use as bounding box] (0,-.8) rectangle (.25,4.8);
  \node[overlay,anchor=base west,inner sep=0pt] at (-4.2,2.8) {$\inserted{\mathrm{Pic}(\overset{\handwritten{\sim}}{X})=}$};
  \draw[overlay] (0,-.8) .. controls (-.15,1.4) and (.35,1.6) .. (1.0,1.7)
    .. controls (2.6,1.8) and (3.6,2.0) .. (4.0,2.3)
    .. controls (4.6,2.6) and (4.8,3.2) .. (4.8,3.8);
}
  \mathrm{Pic}(\overset{\handwritten{\sim}}{X}/\overset{\handwritten{\sim}}{Y})\handwritten{\longrightarrow}
  \mathrm{Pic}(\overset{\handwritten{\sim}}{X}_{\handwritten{\tilde\eta}}/\handwritten{\tilde\eta})\quad ,
\]
where the $\handwritten{\sim}$ denote the effect of strict localization with respect to a geometric point $\bar y$ over $y$. \deleted{When} Given that \deleted{T}\handwritten{t}he preceding morphism factors through the subgroup $\mathrm{Pic}(\overset{\handwritten{\sim}}{X}_{\handwritten{\eta}})$, and that $\mathrm{Pic}(\overset{\handwritten{\sim}}{X})\handwritten{\longrightarrow}\mathrm{Pic}(\overset{\handwritten{\sim}}{X}_{\handwritten{\tilde\eta}})$ is surjective since $\overset{\handwritten{\sim}}{X}$ is regular, we find
\[
  \deleted{\underline{\unsure{E|F}}_{\bar y}}\qquad
  \underline F_{\bar{\deleted{y}\handwritten{y}}^{\deleted{\unsure{*}}}}=
  \mathrm{Pic}(\overset{\handwritten{\sim}}{X}_{\handwritten{\eta}}/\handwritten{\tilde\eta})\,/\,\mathrm{Pic}(\overset{\handwritten{\sim}}{X}_{\handwritten{\tilde\eta}}) \tag*{(4.11)}
\]
\[
  \underline E_{\bar y}\deleted{\unsure{\delta_y|y}}=
  \operatorname{Ker}\bigl(\mathrm{Pic}(\overset{\handwritten{\sim}}{X})\handwritten{\longrightarrow}\mathrm{Pic}(\overset{\handwritten{\sim}}{X}_{\handwritten{\tilde\eta}})\bigr)\quad . \tag*{(4.12)}
\]
Given that $\mathrm{Br}(\handwritten{\tilde\eta})=0$ if $k(y)$ is perfect, and that it is a $p(y)$-torsion group in any case, where $p(y)$ is the chara\deleted{tc}\handwritten{ct}eristic exponent of $k(y)$, we find that $\underline F_{\bar y}$ is a $p(y)$-torsion group, zero if $k(y)$ is perfect, Denoting by $\handwritten{\delta}_{y},\handwritten{\delta'}_{y},\handwritten{\delta''}_{y}$ the integers analogous to $\handwritten{\delta},\handwritten{\delta'},\handwritten{\delta''}$ for the localized situation $\overset{\handwritten{\sim}}{X}\handwritten{\longrightarrow}\overset{\handwritten{\sim}}{Y}$, an argument already \deleted{sig} made above shows us that $\underline F_{\bar y}$ \deleted{is also zero} is also zero if $\handwritten{\delta}_{y}=1$, for example if $\overset{\handwritten{\sim}}{X}$ admits \unsure{a} section over $\overset{\handwritten{\sim}}{Y}$. We shall assume for simplicity, in what follows, that
% archive sheet number: 9
\clearpage
